\section{An integral tableau basis for the exact plethystic maps}
\label{sec:plethysm-integral-tableau-basis}

All constructions in this section are over an arbitrary commutative ring
$R$ with identity. In particular, neither an integer nor a determinant
is inverted. A partition is a finite sequence
$\lambda=(\lambda_1\geq\cdots\geq\lambda_\ell>0)$.
Write $d=|\lambda|$, $t=\lambda_1$, and
$h_j=\#\{r:\lambda_r\geq j\}$ for its column lengths.
The empty partition is allowed, with $d=t=\ell=0$.

\subsection{The intrinsic image and its polynomial realization}

For a finite free $R$-module $E$ define
\begin{align}
 C_\lambda(E)&=\bigotimes_{j=1}^{t}\bigwedge^{h_j}E,
 &P_\lambda(E)&=\bigotimes_{r=1}^{\ell}\operatorname{Sym}^{\lambda_r}E.
 \label{eq:plethysm-column-row-modules}
\end{align}
An empty tensor product in this section is $R$.
For each column, send an exterior product to the alternating tensor
\[
 v_1\wedge\cdots\wedge v_h\longmapsto
 \sum_{\sigma\in\mathfrak S_h}\operatorname{sgn}(\sigma)
 v_{\sigma(1)}\otimes\cdots\otimes v_{\sigma(h)}.
\]
This formula is multilinear and alternating over $R$: when two inputs
are equal, transposing their labels pairs equal tensors with opposite
coefficients. It therefore defines a map out of the exterior power,
also in characteristic two. Tensor these maps in the column order
$1,\ldots,t$, relabel their tensor positions by the boxes $(r,j)$,
reorder the positions into row order, and take the symmetric quotient
in each row. Denote the resulting map by
\begin{equation}
 \delta_{\lambda,E}:C_\lambda(E)\longrightarrow P_\lambda(E),
 \qquad
 \mathsf S_\lambda(E)=\operatorname{im}\delta_{\lambda,E}.
 \label{eq:plethysm-schur-image}
\end{equation}
Equation~\eqref{eq:plethysm-schur-image} is the Schur module convention
used throughout this section. There is no additional row
symmetrizer, factorial, or Young-idempotent normalization.

Let $e_1,\ldots,e_n$ be an ordered basis of $E$.
The module $\operatorname{Sym}^a(E)$ has the basis
$e_1^{\alpha_1}\cdots e_n^{\alpha_n}$ with
$\sum_i\alpha_i=a$. Indeed, the tensor-word basis is partitioned into
permutation orbits, and the defining symmetry relations identify the
words in each orbit; the free module on these orbits gives an inverse
to this quotient description. The module $\bigwedge^a E$ has the basis
$e_{i_1}\wedge\cdots\wedge e_{i_a}$ for
$i_1<\cdots<i_a$. To see this without an assumption on $R$, sort each
word with its permutation sign and send a word with repeated indices
to zero. This alternating multilinear rule gives an inverse to the
map from the free module on the displayed increasing words.

Introduce independent variables $X_{r,i}$ for
$1\leq r\leq\ell$ and $1\leq i\leq n$.
The row monomial bases identify $P_\lambda(E)$ with the component of
$R[X_{r,i}]$ having degree $\lambda_r$ in the variables with first
index $r$. For an ordered list $I=(i_1,\ldots,i_a)$, where
$a\leq\ell$, put
\[
 [I]_a=\det(X_{r,i_s})_{1\leq r,s\leq a},
 \qquad [\varnothing]_0=1.
\]
Repeated entries give zero, and permuting the list multiplies this
polynomial by the permutation sign.
For a filling $T$ of the diagram of $\lambda$, write $I_j(T)$ for its
$j$th column, read from top to bottom, and set
\begin{equation}
 b_T=\prod_{j=1}^{t}[I_j(T)]_{h_j}.
 \label{eq:plethysm-flag-polynomial}
\end{equation}
Expanding the alternating tensors in the definition of $\delta$
shows term by term that $b_T$ is the image of the column exterior
tensor indexed by $T$. Consequently the $b_T$ with strictly increasing
columns generate $\mathsf S_\lambda(E)$. A filling is called
\emph{semistandard} when its columns strictly increase and its rows
weakly increase, using the alphabet $1,\ldots,n$.

\subsection{An alternating identity and an integral straightening proof}

The following identity is proved in the polynomial ring, so it
remains valid after every specialization of its coefficients.

\begin{lemma}[The tail--head shuffle identity]
\label{lem:plethysm-shuffle}
Let $\ell\geq a\geq b\geq k\geq1$ and let $m=a-k+1$.
Let $A$ be a list of column indices of length $k-1$,
let $U=(u_1,\ldots,u_{a+1})$,
and let $B$ be a list of length $b-k$. For a subset
$S=\{s_1<\cdots<s_m\}\subseteq\{1,\ldots,a+1\}$, let $U_S$
be the sublist in the original position order; the complementary
sublist has the same convention. Set
\[
 \epsilon(S)=(-1)^{\sum_{v=1}^{m}(s_v-v)}.
\]
Then
\begin{equation}
 \sum_{\substack{S\subseteq\{1,\ldots,a+1\}\\|S|=m}}
 \epsilon(S)[A,U_S]_a[U_{S^c},B]_b=0.
 \label{eq:plethysm-shuffle}
\end{equation}
Here the minors use initial row segments of the same matrix.
\end{lemma}

\begin{proof}
First let every entry of the lists denote a vector of $R^a$, and
let the second determinant use its projection onto the first $b$
coordinates. Keep $A$ and $B$ fixed. The left side is a multilinear
function of the $a+1$ vectors of $U$. It is alternating, as follows.
Suppose two of these vectors have equal values. A term putting both
in the first determinant, or both in the second determinant, is
zero. Pair each remaining subset with the subset obtained by
exchanging those two positions between the subset and its complement.
The two determinant products in such a pair agree up to the signs
needed to restore the position orders inside the two sublists.
The shuffle sign is the sign of the permutation that lists $S$
first and $S^c$ second. Including the two internal reordering signs,
the exchange is a single transposition of the full list. Thus the
two paired terms have opposite coefficients and cancel. This proves
vanishing on equal inputs by an explicit pairing, without deducing
it from skew symmetry or dividing by two.

Expand each of the $a+1$ vectors in the $a$ standard basis vectors
of $R^a$. Multilinearity expresses the value as a sum of values on
standard basis vectors. Each of these tuples repeats a vector, by
the pigeonhole principle, so each value is zero by the preceding
argument. The multilinear function therefore vanishes identically.
Apply this result over the polynomial ring to the column vectors
$(X_{1,i},\ldots,X_{a,i})^{\mathsf t}$. Its projection onto the first
$b$ coordinates is precisely the column occurring in the second
flag minor. This gives~\eqref{eq:plethysm-shuffle} with its stated
signs and row sets.
\end{proof}

\begin{theorem}[Integral semistandard basis]
\label{thm:plethysm-integral-basis}
For every finite free $R$-module $E$ with its specified ordered basis,
the elements $b_T$ indexed by semistandard tableaux of shape $\lambda$
form an $R$-basis of $\mathsf S_\lambda(E)$.
The image in~\eqref{eq:plethysm-schur-image} is a direct summand as an
$R$-module of $P_\lambda(E)$; the splitting need not be natural.
For every ring homomorphism $R\to R'$, the natural base-change map
\begin{equation}
 R'\otimes_R\mathsf S_\lambda(E)
 \longrightarrow
 \mathsf S_\lambda(R'\otimes_R E)
 \label{eq:plethysm-basechange}
\end{equation}
is an isomorphism. No flatness hypothesis is needed.
\end{theorem}

\begin{proof}
If $\lambda$ is empty, all three claims concern the identity map on
$R$ and follow directly. If $h_1>n$, the first exterior factor is
zero, and there are no strictly increasing first columns. Both the
image and the claimed basis are empty, giving the zero module.
We may therefore suppose that $\lambda$ is nonempty and $h_1\leq n$.

First prove spanning. Order all strictly column-increasing fillings
by lexicographic order of the word obtained by reading their columns
from left to right, and each column from top to bottom. This is a
finite totally ordered set. A filling $T$ that is not semistandard
has adjacent columns
\[
 I=(i_1<\cdots<i_a),\qquad J=(j_1<\cdots<j_b),
 \qquad a\geq b,
\]
and a row $k\leq b$ at which $i_k>j_k$.
Apply Lemma~\ref{lem:plethysm-shuffle} with
\[
 A=(i_1,\ldots,i_{k-1}),\quad
 U=(i_k,\ldots,i_a,j_1,\ldots,j_k),\quad
 B=(j_{k+1},\ldots,j_b).
\]
The subset $S_0=\{1,\ldots,a-k+1\}$ has sign $+1$ and gives
exactly $[I]_a[J]_b$. Move all the other summands to the other side
of~\eqref{eq:plethysm-shuffle}, and multiply by the unchanged minors
in the remaining columns of $T$.

Every member of $(i_k,\ldots,i_a)$ is larger than every member of
$(j_1,\ldots,j_k)$. Thus the entries of $U$ are distinct. A subset
$S\ne S_0$ replaces at least one entry of the left-column tail by
a strictly smaller entry from the right-column head. If a resulting
column has a repeated entry, its determinant is zero and the term
is discarded. In any other term, sort the two columns numerically
and multiply by the two sorting signs. The fixed left-column prefix
is unchanged as a multiset; all inserted values are smaller than
all removed values. At the first place where the sorted old and
new left columns differ, the new column therefore has the smaller
entry. All preceding columns of the entire filling are unchanged.
The resulting filling is strictly smaller than $T$ in the specified
total order. The coefficient of the original term was exactly one,
so this replacement uses only integer addition and subtraction.
Induction on the finite total order proves that every generator
$b_T$ is an integral linear combination of the semistandard $b_T$.
This proof works after passage to $R$ as well.

For independence, order the variables by
\[
 X_{1,1}>X_{1,2}>\cdots>X_{1,n}>
 X_{2,1}>\cdots>X_{2,n}>\cdots>X_{\ell,n},
\]
and use lexicographic monomial order. If
$i_1<\cdots<i_a$, the leading monomial of $[I]_a$ is
$X_{1,i_1}\cdots X_{a,i_a}$, with coefficient $+1$.
Indeed, among determinant permutations the largest monomial must
choose the smallest available column index for its first row, then
the smallest remaining index for its second row, and so on.
Monomial order is preserved by multiplication, so
\begin{equation}
 \operatorname{in}(b_T)=
 \prod_{(r,j)\in\lambda}X_{r,T(r,j)},
 \qquad
 \text{with coefficient }+1.
 \label{eq:plethysm-leading-monomial}
\end{equation}
For a semistandard filling, this monomial determines each row's
multiset of entries. A weakly increasing row is uniquely recovered
from that multiset. Hence distinct semistandard tableaux have
distinct leading monomials. In a nonzero linear combination over
$R$, choose the largest leading monomial among the tableaux whose
coefficients are nonzero. No other summand has a monomial that
large. Its coefficient in the sum is the chosen nonzero coefficient
times $1$, and cannot vanish. This proves independence even when
$R$ has zero divisors.

For completeness, this monicity also produces a splitting. Work
first over $\mathbb Z$ and list the semistandard polynomials in
decreasing order of their distinct leading monomials. Given any
polynomial in the finite row-degree component $P_\lambda(\mathbb Z^n)$,
process this list in order. At each stage, extract the coefficient
of the current leading monomial, record it, and subtract that
coefficient times the corresponding $b_T$. Later subtractions
cannot change a leading monomial already processed, since all
their monomials are smaller. Return the sum of the recorded
multiples of $b_T$. This is a $\mathbb Z$-linear map into their
span, and it is the identity on each $b_T$. It is therefore a
retraction of the inclusion of that span. All its matrix entries
are integers, so base change gives the stated retraction over
every $R$.

Finally the map~\eqref{eq:plethysm-basechange} sends
$1\otimes b_T$ to the polynomial $b_T$ with coefficients in $R'$.
The basis just proved on both sides shows that it is an
isomorphism. The splitting explains in particular why taking this
image commutes with arbitrary, possibly nonflat, base change.
\end{proof}

\subsection{Functoriality, all matrix entries, and the exact diagonal}

\begin{proposition}[Functorial maps without denominators]
\label{prop:plethysm-functorial}
Every $R$-linear map $f:E\to F$ between finite free modules induces
a map $\mathsf S_\lambda(f):\mathsf S_\lambda(E)\to
\mathsf S_\lambda(F)$. These maps preserve identities and
composition and commute with arbitrary base change. In the tableau
bases, every matrix entry is an integer polynomial, homogeneous of
degree $|\lambda|$ in the entries of $f$.
\end{proposition}

\begin{proof}
The alternating tensor map commutes with $f^{\otimes h}$ term by
term, tensor reordering commutes with applying $f$ in every slot,
and symmetric quotient maps are natural. Thus the square
\[
 \begin{array}{ccc}
 C_\lambda(E)&\xrightarrow{\delta_{\lambda,E}}&P_\lambda(E)\\
 \big\downarrow C_\lambda(f)&&\big\downarrow P_\lambda(f)\\
 C_\lambda(F)&\xrightarrow{\delta_{\lambda,F}}&P_\lambda(F)
 \end{array}
\]
commutes. Restricting its right arrow to the images defines the
claimed map. The identities and composition laws hold on the
ambient row-symmetric modules and hence on their submodules.
All constructions use tensor, exterior, and symmetric operations
given by the displayed integral formulas. Together with
Theorem~\ref{thm:plethysm-integral-basis}, this proves the base-change
statement.

Here is a formula that also proves the claim about every matrix
entry. Choose bases $e_i$ and $f_j$, and write a matrix $M$ for the
map, so that $f(e_i)=\sum_jM_{j,i}f_j$. For a strictly increasing
column $I$ of length $a$, direct multilinear expansion and signed
sorting in the exterior power give
\begin{equation}
 (\bigwedge^a f)(e_{i_1}\wedge\cdots\wedge e_{i_a})
 =\sum_{\substack{J\text{ increasing}\\|J|=a}}
       \det(M_{J,I})\,f_{j_1}\wedge\cdots\wedge f_{j_a}.
 \label{eq:plethysm-exterior-matrix}
\end{equation}
Terms with repeated row indices vanish. For each distinct row
set the terms with the possible row orders sum, with their sorting
signs, to the determinant in~\eqref{eq:plethysm-exterior-matrix}.
For the columns $I_c(T)$ of a tableau $T$, take the product of
these expansions and then apply $\delta_{\lambda,F}$. This gives
\begin{equation}
 \mathsf S_\lambda(f)b_T
 =\sum_{J_1,\ldots,J_t}
       \left(\prod_{c=1}^t\det(M_{J_c,I_c(T)})\right)
       b_{(J_1,\ldots,J_t)}.
 \label{eq:plethysm-general-matrix}
\end{equation}
Apply the terminating integer straightening procedure already
proved to each filling on the right. Every product of minors in
its coefficient has total degree $\sum_c h_c=|\lambda|$.
Straightening multiplies by integers and adds these products, so
each final tableau coefficient is an integer homogeneous
polynomial of that degree. This is a finite formula: there are
$\prod_c\binom{\operatorname{rank}F}{h_c}$ choices in the sum,
and straightening strictly decreases an order on the finite set
of column-increasing fillings. No complexity bound beyond this
finite procedure is asserted here.
\end{proof}

\begin{theorem}[The exact signed diagonal map]
\label{thm:plethysm-diagonal}
Let $E$ and $F$ have ordered bases of the same cardinality $n$,
and let $D:E\to F$ have the specified matrix
$\operatorname{diag}(d_1,\ldots,d_n)$, with arbitrary $d_i\in R$.
For a semistandard tableau $T$, let $c_i(T)$ count its entries
equal to $i$, and put
\[
 w_T=\prod_{i=1}^n d_i^{c_i(T)}.
\]
With the same tableau order on the two bases, the exact matrix of
$\mathsf S_\lambda(D)$ is $\operatorname{diag}(w_T)_T$.
In particular, there are the following specified maps and modules:
\begin{align}
 \ker\mathsf S_\lambda(D)
   &=\bigoplus_T\operatorname{Ann}_R(w_T)\,b_T(E),\\
 \operatorname{im}\mathsf S_\lambda(D)
   &=\bigoplus_T(w_T)\,b_T(F),\\
 \operatorname{coker}\mathsf S_\lambda(D)
   &\cong\bigoplus_T R/(w_T).
 \label{eq:plethysm-diagonal-modules}
\end{align}
The last isomorphism sends the class of
$\sum_T a_Tb_T(F)$ to $(a_T\bmod (w_T))_T$.
\end{theorem}

\begin{proof}
Apply $D$ in the column exterior tensor defining $b_T$.
Each occurrence of $e_i$ contributes exactly the scalar $d_i$,
and the product of these scalars is $w_T$. The commutative
square in Proposition~\ref{prop:plethysm-functorial} therefore
gives $\mathsf S_\lambda(D)b_T(E)=w_Tb_T(F)$ before any
straightening. The tableau basis theorem proves the asserted
matrix identity. A vector is in its kernel precisely when every
coordinate $a_T$ satisfies $w_Ta_T=0$, which gives the first
formula. Each target coordinate of an image is independently an
element of $(w_T)$, giving the second. The displayed residue map
is surjective and has exactly this image as its kernel, proving
the third. This also treats $d_i=0$ and nonzero zero divisors.
Signed elements $d_i$ have not been replaced by associates or
absolute values.
\end{proof}

\subsection{Iteration and the plethystic exponent matrix}

Let $\mathcal T_\lambda(n)$ be the tableau set in
Theorem~\ref{thm:plethysm-integral-basis}, ordered by its
column-reading words. Put $N=|\mathcal T_\lambda(n)|$.
Regard this ordered set as the alphabet for semistandard tableaux
$U$ of a second partition $\mu$. Write $c_T(U)$ for the number
of entries of $U$ equal to the alphabet letter $T$, and define
\begin{equation}
 \Gamma_i(U)=\sum_{T\in\mathcal T_\lambda(n)}c_T(U)c_i(T).
 \label{eq:plethysm-nested-exponents}
\end{equation}
Here the plethysm is the composed functor
$\mathsf S_\mu\circ\mathsf S_\lambda$, with this integral
Schur-image convention at both stages.

\begin{corollary}[Nested Schur maps over the original ring]
\label{cor:plethysm-nested-diagonal}
The outer tableau elements $b_U$ form an $R$-basis of
$\mathsf S_\mu(\mathsf S_\lambda(E))$.
For the specified diagonal map $D$ of
Theorem~\ref{thm:plethysm-diagonal}, the exact nested map is
\begin{equation}
 \mathsf S_\mu(\mathsf S_\lambda(D))b_U(E)
   =\left(\prod_{i=1}^n d_i^{\Gamma_i(U)}\right)b_U(F).
 \label{eq:plethysm-nested-diagonal}
\end{equation}
Every exponent is a nonnegative integer, and
$\sum_i\Gamma_i(U)=|\mu|\,|\lambda|$.
The kernel, image, and cokernel are given by
\eqref{eq:plethysm-diagonal-modules} with the index $U$ and
scalar $\prod_i d_i^{\Gamma_i(U)}$.
For a general matrix $M$, every matrix entry of the composed
functor is an integer polynomial homogeneous of degree
$|\mu|\,|\lambda|$ in $M$.
\end{corollary}

\begin{proof}
The inner module is finite free with the explicitly ordered basis
$(b_T)_{T\in\mathcal T_\lambda(n)}$ by the basis theorem.
Apply the same basis theorem to that free module and $\mu$.
The inner map is diagonal, with entries
$w_T=\prod_i d_i^{c_i(T)}$. Apply the diagonal theorem to this
map and the outer shape $\mu$. Its scalar on $b_U$ is exactly
\[
 \prod_T w_T^{c_T(U)}
   =\prod_T\prod_i d_i^{c_T(U)c_i(T)}
   =\prod_i d_i^{\Gamma_i(U)}.
\]
The number of letters of an inner tableau is $|\lambda|$, so
$\sum_i c_i(T)=|\lambda|$; the number of letters of an outer
tableau is $|\mu|$, so $\sum_T c_T(U)=|\mu|$. Summing
\eqref{eq:plethysm-nested-exponents} proves the degree identity.
The coordinate kernel and quotient formulas follow from the
diagonal theorem, with no additional ring hypotheses.
For a general $M$, apply the integer homogeneous matrix formula
first to $\lambda$, and then substitute those degree-$|\lambda|$
entries into the degree-$|\mu|$ formula for the outer functor.
This proves the last assertion. Empty shapes and zero inner
rank follow directly from the conventions for empty tensor
products and the absence of nonempty tableaux on an empty
alphabet.
\end{proof}

The construction above retains the entire composed functor and
its tableau matrix. It does not require a decomposition into
irreducible Schur summands. The proof gives ordinary integral
Schur operations on the specified linear maps; identifying them
with a particular nonstandard quantum-group construction requires
the corresponding specified representation map.
For a primary comparison with the flag-minor standard monomial
model, see D.~C.~Lax, \emph{Accessible Proof of Standard Monomial
Basis for Coordinatization of Schubert Sets of Flags},
\href{https://arxiv.org/abs/1508.02744}{arXiv:1508.02744},
Sections~3--4. All basis, straightening, functoriality, and
diagonal assertions used here were proved above over $R$.
